\documentclass[11pt]{article}
% ===========================================================================
%  Shared preamble for all MPM2D worksheets — input right after \documentclass.
%  (Files beginning with "_" are skipped by mpm2d-worksheet-publish.mjs.)
% ===========================================================================
\usepackage[letterpaper,top=2.2cm,bottom=1.9cm,left=1.6cm,right=1.6cm,headheight=28pt,footskip=24pt]{geometry}
\usepackage{amsmath,amssymb}
\usepackage[T1]{fontenc}
\usepackage[utf8]{inputenc}
\usepackage{lmodern}
\usepackage{xcolor}
\usepackage[most]{tcolorbox}
\usepackage{enumitem}
\usepackage{multicol}
\usepackage{fancyhdr}
\usepackage{tikz}
\usetikzlibrary{calc}
\usepackage{pgfplots}
\pgfplotsset{compat=1.18}

\definecolor{exblue}{HTML}{4A90E2}
\definecolor{exbg}{HTML}{E6F3FF}
\definecolor{qorange}{HTML}{E69138}
\definecolor{qbg}{HTML}{FFF7CC}
\definecolor{keygray}{HTML}{475569}
\definecolor{keybg}{HTML}{F1F5F9}
\definecolor{iagreen}{HTML}{14653B}
\definecolor{titleblue}{HTML}{1D4ED8}
% Rotating palette so each worked example gets its own colour.
\definecolor{ex0}{HTML}{2563A0}\definecolor{ex1}{HTML}{3B7D3B}\definecolor{ex2}{HTML}{8A5A00}
\definecolor{ex3}{HTML}{A3327A}\definecolor{ex4}{HTML}{6D28A3}\definecolor{ex5}{HTML}{B08900}
\definecolor{ex6}{HTML}{0E7490}\definecolor{ex7}{HTML}{9A3412}\definecolor{ex8}{HTML}{15803D}

\pagestyle{fancy}
\fancyhf{}
\fancyhead[L]{\footnotesize NAME: \underline{\hspace{3.4cm}}}
\fancyhead[C]{\textbf{Integration Academy}\\[-2pt]{\scriptsize Math Simplified \textperiodcentered{} Success Amplified}}
\fancyhead[R]{\footnotesize MPM2D}
\fancyfoot[L]{\scriptsize dr.merie@IntegrationAcademy.ca}
\fancyfoot[C]{\scriptsize\thepage}
\fancyfoot[R]{\scriptsize IntegrationAcademy.ca}
\renewcommand{\headrulewidth}{1.2pt}
\renewcommand{\footrulewidth}{0.4pt}
\setlength{\parindent}{0pt}
\setlength{\parskip}{4pt}

% Examples are NOT breakable, so each example + its solution stays on one page.
\newtcolorbox{example}[2]{colframe=#1,colback=#1!7,coltitle=white,colbacktitle=#1,fonttitle=\bfseries,title={#2},arc=3pt,boxrule=1pt}
\newtcolorbox{qbox}[1]{colframe=qorange,colback=qbg,coltitle=white,colbacktitle=qorange,fonttitle=\bfseries,title={#1},arc=3pt,breakable,boxrule=1pt}
\newtcolorbox{keybox}{colframe=keygray,colback=keybg,coltitle=white,colbacktitle=keygray,fonttitle=\bfseries,title={$\checkmark$\ Answer Key},arc=3pt,breakable,boxrule=1pt}
\newtcolorbox{recapbox}{colframe=keygray,colback=keybg,coltitle=white,colbacktitle=keygray,fonttitle=\bfseries,title={Question Recap},arc=3pt,breakable,boxrule=1pt}
\newtcolorbox{ideabox}{colframe=titleblue,colback=titleblue!6,coltitle=white,colbacktitle=titleblue,fonttitle=\bfseries,title={The Big Ideas},arc=3pt,breakable,boxrule=1.2pt}
\newtcolorbox{learnbox}{colframe=iagreen,colback=iagreen!5,coltitle=white,colbacktitle=iagreen,fonttitle=\bfseries,title={Core Concepts},arc=3pt,breakable,boxrule=1.2pt}
\newcommand{\lh}[1]{\par\smallskip{\bfseries\color{iagreen}#1}\par\nobreak\smallskip}

\newcommand{\soln}{\par\smallskip\textit{\textbf{Solution.}}\ }
\newcommand{\workspace}[1]{\par\vspace{3pt}\fcolorbox{black!30}{white}{\begin{minipage}[t][#1][t]{\dimexpr\linewidth-2\fboxsep\relax}\ \end{minipage}}\par}
\newcommand{\sech}[1]{\par\vspace{8pt}{\Large\bfseries\color{iagreen}#1}\par\nobreak\vspace{2pt}{\color{black!20}\hrule height1pt}\vspace{6pt}}

% A compact coordinate plot sized for an example box: \eplot{xmin}{xmax}{ymin}{ymax}{body}
\newcommand{\eplot}[5]{\begin{center}\begin{tikzpicture}
\begin{axis}[width=6.8cm,height=4.4cm,axis lines=middle,xlabel={\tiny$x$},ylabel={\tiny$y$},
grid=both,grid style={gray!16},xmin=#1,xmax=#2,ymin=#3,ymax=#4,
xtick distance=2,ytick distance=2,tick label style={font=\tiny},axis line style={-},samples=120]
#5
\end{axis}\end{tikzpicture}\end{center}}

% Same as \eplot but with its own tick spacing: \eplotx{xmin}{xmax}{ymin}{ymax}{xtick}{ytick}{body}
\newcommand{\eplotx}[7]{\begin{center}\begin{tikzpicture}
\begin{axis}[width=8.4cm,height=5.4cm,axis lines=middle,xlabel={\tiny$x$},ylabel={\tiny$y$},
grid=both,grid style={gray!16},xmin=#1,xmax=#2,ymin=#3,ymax=#4,
xtick distance=#5,ytick distance=#6,tick label style={font=\tiny},axis line style={-},samples=120]
#7
\end{axis}\end{tikzpicture}\end{center}}

% A sine-curve plot (degrees on x, 0..360): \splot{ymin}{ymax}{body}
\newcommand{\splot}[3]{\begin{center}\begin{tikzpicture}
\begin{axis}[width=8.6cm,height=4.6cm,axis lines=middle,xlabel={\tiny$x\,(^\circ)$},ylabel={\tiny$y$},
grid=both,grid style={gray!16},xmin=0,xmax=360,ymin=#1,ymax=#2,
xtick distance=90,ytick distance=1,tick label style={font=\tiny},axis line style={-},samples=180]
#3
\end{axis}\end{tikzpicture}\end{center}}

% A small coordinate plot: \plot{xmin}{xmax}{ymin}{ymax}{body}
\newcommand{\plot}[5]{\begin{center}\begin{tikzpicture}
\begin{axis}[width=8.4cm,height=6cm,axis lines=middle,xlabel={\scriptsize$x$},ylabel={\scriptsize$y$},
grid=both,grid style={gray!16},xmin=#1,xmax=#2,ymin=#3,ymax=#4,
xtick distance=2,ytick distance=2,tick label style={font=\tiny},axis line style={-}]
#5
\end{axis}\end{tikzpicture}\end{center}}

% Right triangle (angle theta at bottom-left, right angle at bottom-right):
%   \rtri{drawBase}{drawHeight}{oppLabel}{adjLabel}{hypLabel}
\newcommand{\rtri}[5]{\begin{center}\begin{tikzpicture}[scale=0.85]
\coordinate (A) at (0,0);\coordinate (B) at (#1,0);\coordinate (C) at (#1,#2);
\draw[thick,exblue] (A)--(B)--(C)--cycle;
\draw[black] ($(B)+(-0.32,0)$)--($(B)+(-0.32,0.32)$)--($(B)+(0,0.32)$);
\node[below] at ($(A)!0.5!(B)$){#4};
\node[right] at ($(B)!0.5!(C)$){#3};
\node[above left] at ($(A)!0.5!(C)$){#5};
\node at (0.7,0.22){$\theta$};
\end{tikzpicture}\end{center}}

% General (oblique) triangle with vertex labels and opposite-side labels:
%   \gtri{Alabel}{Blabel}{Clabel}{aSide}{bSide}{cSide}
\newcommand{\gtri}[6]{\begin{center}\begin{tikzpicture}[scale=0.85]
\coordinate (A) at (0,0);\coordinate (B) at (5,0);\coordinate (C) at (1.6,3);
\draw[thick,exblue] (A)--(B)--(C)--cycle;
\node[below left] at (A){#1};\node[below right] at (B){#2};\node[above] at (C){#3};
\node[below] at ($(A)!0.5!(B)$){#6};
\node[right] at ($(B)!0.5!(C)$){#4};
\node[left] at ($(A)!0.5!(C)$){#5};
\end{tikzpicture}\end{center}}

\fancyhead[R]{\footnotesize MHF4U \textperiodcentered{} 6.3}
\begin{document}

\begin{center}
{\LARGE\bfseries\color{titleblue}Worksheet 6.3 --- Proving Trigonometric Identities}\\[2pt]
{\small Grade 12 Advanced Functions (MHF4U) \textperiodcentered{} Unit 6: Trigonometric Identities and Equations}
\end{center}

\begin{ideabox}
Transform one side using the Pythagorean, quotient, and reciprocal identities until it matches the other.
\begin{itemize}[leftmargin=*,itemsep=2pt]
  \item $\sin^2\theta+\cos^2\theta=1$; $1+\tan^2\theta=\sec^2\theta$; $1+\cot^2\theta=\csc^2\theta$.
  \item $\tan\theta=\tfrac{\sin\theta}{\cos\theta}$, $\cot\theta=\tfrac{\cos\theta}{\sin\theta}$.
  \item Start with the harder side; rewrite in sines and cosines.
\end{itemize}
\end{ideabox}

\begin{learnbox}
\lh{What a proof shows}An \textbf{identity} holds for all valid angles; a proof transforms one side into the other using known identities and algebra.
\lh{A steady strategy}Start from the more complex side, rewrite everything in sines and cosines, and lean on the Pythagorean and compound identities.
\lh{Not an equation}You never move terms across the equals sign — each step simply rewrites one side until the two match.
\end{learnbox}

\sech{Worked Examples}

\begin{example}{ex0}{Example 1 --- Pythagorean}
Prove $\dfrac{\sin^2 x}{1-\cos x}=1+\cos x$.\soln $\sin^2 x=1-\cos^2 x$ (the sum below is always $1$): $\dfrac{(1-\cos x)(1+\cos x)}{1-\cos x}=1+\cos x$. $\checkmark$\begin{center}\begin{tikzpicture}\begin{axis}[width=9.2cm,height=4.6cm,axis lines=middle,xlabel={\small$x$},ylabel={\small$y$},xmin=-0.4,xmax=6.9,ymin=-0.4,ymax=2,samples=220,xtick={1.5708,3.1416,4.7124,6.2832},xticklabels={$\frac{\pi}{2}$,$\pi$,$\frac{3\pi}{2}$,$2\pi$}]\addplot[exblue,very thick,domain=0:6.4]{(sin(deg(x)))^2+(cos(deg(x)))^2};\end{axis}\end{tikzpicture}\end{center}
\end{example}

\begin{example}{ex1}{Example 2 --- Quotient}
Prove $\tan x\cos x=\sin x$.\soln $\dfrac{\sin x}{\cos x}\cdot\cos x=\sin x$. $\checkmark$
\end{example}

\begin{example}{ex2}{Example 3 --- Reciprocal ratios}
Prove $\dfrac{\tan x}{\sin x}=\sec x$.\soln $\dfrac{\sin x/\cos x}{\sin x}=\dfrac1{\cos x}=\sec x$. $\checkmark$
\end{example}

\begin{example}{ex3}{Example 4 --- Build to a known one}
Prove $\cot x\sec x=\csc x$.\soln $\dfrac{\cos x}{\sin x}\cdot\dfrac1{\cos x}=\dfrac1{\sin x}=\csc x$. $\checkmark$
\end{example}

\begin{example}{ex4}{Example 5 --- Pythagorean form}
Prove $\cos^2 x\,(1+\tan^2 x)=1$.\soln $1+\tan^2 x=\sec^2 x$, so $\cos^2 x\cdot\dfrac1{\cos^2 x}=1$. $\checkmark$
\end{example}

\begin{example}{ex5}{Example 6 --- Difference of squares}
Prove $(1-\sin x)(1+\sin x)=\cos^2 x$.\soln $1-\sin^2 x=\cos^2 x$. $\checkmark$
\end{example}

\begin{example}{ex6}{Example 7 --- Quotient}
Prove $\cot x\sin x=\cos x$.\soln $\dfrac{\cos x}{\sin x}\cdot\sin x=\cos x$. $\checkmark$
\end{example}

\begin{example}{ex7}{Example 8 --- Quotient}
Prove $\dfrac{\cos x}{\cot x}=\sin x$.\soln $\cos x\cdot\dfrac{\sin x}{\cos x}=\sin x$. $\checkmark$
\end{example}

\begin{example}{ex8}{Example 9 --- Combine}
Prove $\sec x-\cos x=\sin x\tan x$.\soln $\dfrac1{\cos x}-\cos x=\dfrac{1-\cos^2 x}{\cos x}=\dfrac{\sin^2 x}{\cos x}=\sin x\cdot\dfrac{\sin x}{\cos x}$. $\checkmark$
\end{example}

\clearpage
\sech{Your Turn --- Show Your Work}
For each question, show all steps clearly in the space provided.

\begin{qbox}{Question 1}
Prove $\cos^2 x-1=-\sin^2 x$.\workspace{2.4cm}
\end{qbox}

\begin{qbox}{Question 2}
Prove $\cot x\tan x=1$.\workspace{2.4cm}
\end{qbox}

\begin{qbox}{Question 3}
Prove $\dfrac{\cot x}{\csc x}=\cos x$.\workspace{2.4cm}
\end{qbox}

\begin{qbox}{Question 4}
Prove $\dfrac{\sec x}{\tan x}=\csc x$.\workspace{2.4cm}
\end{qbox}

\begin{qbox}{Question 5}
Prove $\sin^2 x\,(1+\cot^2 x)=1$.\workspace{2.4cm}
\end{qbox}

\begin{qbox}{Question 6}
Prove $\dfrac{\sin x}{\tan x}=\cos x$.\workspace{2.4cm}
\end{qbox}

\begin{qbox}{Question 7}
Prove $\dfrac{\csc x}{\cot x}=\sec x$.\workspace{2.4cm}
\end{qbox}

\begin{qbox}{Question 8}
Prove $\sin x\sec x=\tan x$.\workspace{2.4cm}
\end{qbox}

\begin{qbox}{Question 9}
Prove $\cos x\csc x=\cot x$.\workspace{2.4cm}
\end{qbox}

\begin{qbox}{Question 10}
Prove $\sec^2 x-\tan^2 x=1$.\workspace{2.4cm}
\end{qbox}

\begin{qbox}{Question 11}
Prove $\dfrac{1-\cos^2 x}{\sin x}=\sin x$.\workspace{2.4cm}
\end{qbox}

\begin{qbox}{Question 12}
Prove $\dfrac{\cos x}{1-\sin^2 x}=\sec x$.\workspace{2.4cm}
\end{qbox}

\begin{qbox}{Question 13 --- Challenge}
Prove $\dfrac{\sin x}{1+\cos x}+\dfrac{1+\cos x}{\sin x}=\dfrac{2}{\sin x}$.\workspace{4cm}
\end{qbox}

\clearpage
\sech{Question Recap \& Answer Key}

\begin{recapbox}
\begin{multicols}{2}
\small
\begin{enumerate}[leftmargin=*,itemsep=3pt]
  \item Prove $\cos^2 x-1=-\sin^2 x$.
  \item Prove $\cot x\tan x=1$.
  \item Prove $\dfrac{\cot x}{\csc x}=\cos x$.
  \item Prove $\dfrac{\sec x}{\tan x}=\csc x$.
  \item Prove $\sin^2 x\,(1+\cot^2 x)=1$.
  \item Prove $\dfrac{\sin x}{\tan x}=\cos x$.
  \item Prove $\dfrac{\csc x}{\cot x}=\sec x$.
  \item Prove $\sin x\sec x=\tan x$.
  \item Prove $\cos x\csc x=\cot x$.
  \item Prove $\sec^2 x-\tan^2 x=1$.
  \item Prove $\dfrac{1-\cos^2 x}{\sin x}=\sin x$.
  \item Prove $\dfrac{\cos x}{1-\sin^2 x}=\sec x$.
  \item Prove $\dfrac{\sin x}{1+\cos x}+\dfrac{1+\cos x}{\sin x}=\dfrac{2}{\sin x}$.
\end{enumerate}
\end{multicols}
\end{recapbox}

\begin{keybox}
\begin{multicols}{2}
\small
\begin{enumerate}[leftmargin=*,itemsep=3pt]
  \item $-(1-\cos^2 x)=-\sin^2 x$
  \item $\tfrac{\cos x}{\sin x}\cdot\tfrac{\sin x}{\cos x}=1$
  \item $\tfrac{\cos x/\sin x}{1/\sin x}=\cos x$
  \item $\tfrac{1/\cos x}{\sin x/\cos x}=\tfrac1{\sin x}$
  \item $\sin^2 x\csc^2 x=1$
  \item $\sin x\cdot\tfrac{\cos x}{\sin x}=\cos x$
  \item $\tfrac{1/\sin x}{\cos x/\sin x}=\tfrac1{\cos x}$
  \item $\sin x\cdot\tfrac1{\cos x}=\tan x$
  \item $\cos x\cdot\tfrac1{\sin x}=\cot x$
  \item from $1+\tan^2 x=\sec^2 x$
  \item $\tfrac{\sin^2 x}{\sin x}=\sin x$
  \item $\tfrac{\cos x}{\cos^2 x}=\tfrac1{\cos x}$
  \item common denom $\to\tfrac{\sin^2 x+(1+\cos x)^2}{\sin x(1+\cos x)}=\tfrac{2}{\sin x}$
\end{enumerate}
\end{multicols}
\end{keybox}

\end{document}
